%%
% BIThesis 研究生学位论文模板 The BIThesis Template for Graduate Thesis
% This file has no copyright assigned and is placed in the Public Domain.
%%

\begin{conclusion}

本文面向黑盒大语言模型输出公平性治理问题，围绕在无法访问模型内部信息的条件下识别偏见，并在不过度破坏原始回复内容的前提下完成有效校正，这两个核心目标展开研究。针对现有方法在黑盒适用性、高置信验证以及细粒度输出修正方面的不足，本文从反事实差异分析出发，构建了覆盖偏见识别、结果验证与定向重写的完整方法框架，并通过系统实验对所提方法的有效性进行了验证。本文的主要成果和贡献如下：

第一，围绕黑盒场景下偏见难以稳定识别的问题，本文构建了一种结合反事实输入比较与二阶段核验的偏见判别方法。该方法以敏感属性替换后的反事实输入为参照，通过比较原始回复与对应反事实回复在语义、情感、毒性和生成倾向等方面的变化情况，对潜在偏见风险进行初步量化。此外，本文针对关键差异命题设计了一个同时考察事实依据与公平风险的反思式验证环节，以减少仅根据表层差异作出判断时可能产生的误检。实验结果显示，本文方法在偏见识别准确率上取得了优于基线方法的表现，消融实验也进一步验证了多维差异建模和反思式验证两个环节对整体识别性能的贡献。该方法为黑盒大语言模型输出偏见的实例级识别提供了更具解释性和稳健性的实现路径。

第二，针对黑盒模型输出后处理中普遍存在的整体改写过重、语义保真不足和回复可用性下降等问题，本文提出了一种基于优先级排序与最小编辑约束的偏见重写方法。该方法包括偏见片段定位、重写优先级排序和最小编辑重写三个阶段。首先，通过句子级与短语级划分，并结合多维评分函数定位需要重点处理的局部偏见片段。然后，引入编辑代价构建重写优先级函数，对候选片段按照风险程度与修改成本进行排序。最后，针对待处理片段生成局部候选替换片段，并通过公平性、语义一致性和编辑代价三项指标顺位选择最优结果。在此基础上按照偏见片段的处理优先级顺序，迭代执行局部重写。实验结果表明，本文方法在语义保持度与偏见去除通过率上均较现有方法有所提升。这表明本文方法能够在尽量少改动原始回复的前提下，更有效地完成黑盒大语言模型输出的偏见缓解。

尽管本文取得了一定成果，但仍然存在一些不足。首先，本文实验主要基于 BOLD 和 RedditBias 数据集展开，虽然覆盖了性别、种族、宗教和政治立场等多类敏感属性，并结合外部评审机制对结果进行了分析，但与真实大规模线上交互场景相比，实验数据的语境复杂度和任务多样性仍然有限。其次，二阶段验证与局部重写流程虽然提升了结果的可信度与可控性，但也引入了额外的推理开销，在更高吞吐量或更低时延的应用场景中仍需进一步优化。

未来研究可从以下几个方向继续展开。一是进一步扩展实验范围，在更多模型、更多任务类型和更复杂真实交互场景中验证方法的稳健性与泛化能力，并探索跨语言、跨文化语境下的偏见识别与修正问题。二是结合更高效的裁判模型、自动化一致性评估和轻量级重写策略，降低双阶段验证与输出重写的计算成本，提升方法在实际系统中的部署效率，推动黑盒大语言模型在公平、可信和可持续部署方向上的深入应用。

\end{conclusion}
